% Part: first-order-logic
% Chapter: axiomatic-deduction
% Section: provability-quantifiers

% verification of properties of provability needed for maximally
% consistent sets in the completeness chapter.

\documentclass[../../../include/open-logic-section]{subfiles}

\begin{document}

\olfileid{fol}{axd}{qpr}

\olsection{\usetoken{S}{derivability} والمكمّمات}

\begin{explain}
  تتطلب مبرهنة الاكتمال أيضًا أن تنتج الاشتقاقات البديهية الحقائق
  المتعلقة بـ$\Proves$ التي نثبتها في هذا القسم.
\end{explain}

\begin{thm}
\ollabel{thm:strong-generalization} إذا كان $c$ !!a{constant} لا يرد
في $\Gamma$ ولا في $!A(x)$، وكان $\Gamma \Proves !A(c)$، فإن $\Gamma
\Proves \lforall[x][!A(x)]$.
\end{thm}

\begin{proof}
% تصحيح برهاني (0123-I1): لا يلزم أن يكون c من مخزون الثوابت المميزة؛ نجدد ثوابت البرهان ثم نبدله بثابت من E.
خذ برهانًا منتهيًا~$\Pi$ لـ~$!A(c)$ من~$\Gamma$، ولتكن~$\Gamma_0$
مجموعة فروضه المستعملة. جدّد ثوابته المميَّزة بواسطة
\olref[qua]{lem:qr-freshening} مع مجموعة المنع~$\{c\}$، ثم اختر
$e\in\mathcal E\setminus\{c\}$ لا يرد في أي صيغة من البرهان المجدَّد ولا يساوي ثابتًا
مميَّزًا على إحدى عقد~\QR{} فيه. بدّل الثابتين~$c$ و~$e$ تقابليًا في
البرهان كله. يظل~$\Gamma_0$ ثابتًا تحت هذا التبديل، لأن كليهما لا يرد
فيه، وتحفظ إعادة التسمية البديهيات وإثبات المقدَّم وصورتي~\QR{} وشروطهما؛
ومن ثم~$\Gamma_0 \Proves !A(e)$.

نضم إلى ذلك الحالة~$!A(e)\lif(\ltrue\lif !A(e))$ من بديهية اللزوم
الأولى، فنحصل بإثبات المقدَّم على
$\Gamma_0\Proves\ltrue\lif !A(e)$. ولأن~$e\in\mathcal E$ ولا يرد في
$\Gamma_0$ ولا~$\ltrue$ ولا~$!A(x)$، تجيز~\QR{} الاستنتاج
$\Gamma_0\Proves\ltrue\lif\lforall[x][!A(x)]$. ومع
$\Gamma_0\Proves\ltrue$ نحصل بإثبات المقدَّم على
$\Gamma_0\Proves\lforall[x][!A(x)]$، ثم تعطينا الرتابة النتيجة المطلوبة
من~$\Gamma$.
\end{proof}

\begin{prop}
\ollabel{prop:provability-quantifiers}
إذا كان $t$ حدًا مغلقًا، فإن:
\begin{tagenumerate}{prvEx,prvAll}
\tagitem{prvEx}{$!A(t) \Proves \lexists[x][!A(x)]$.}{}

\tagitem{prvAll}{$\lforall[x][!A(x)] \Proves !A(t)$.}{}
\end{tagenumerate}
\end{prop}

\begin{proof}
\begin{tagenumerate}{prvEx,prvAll}
\tagitem{prvEx}{من \olref[qua]{ax:q2} ومبرهنة الاستنباط.}{}
\tagitem{prvAll}{من \olref[qua]{ax:q1} ومبرهنة الاستنباط.}{}
\end{tagenumerate}
\end{proof}

\end{document}
